The design and performance of the Butterworth, Chebyshev Type I, Chebyshev Type II, and Elliptic analog bandpass filters were evaluated based on the given laboratory specifications. The system required passband edges at $2000$ and $4000$ rad/sec with a maximum ripple of $1$ dB ($R_p$), and stopband edges at $1500$ and $4500$ rad/sec with a minimum attenuation of $60$ dB ($R_s$). An ideal bandpass filter was also modeled using symbolic expressions to act as a reference for comparison.

As shown in Figure 1, the magnitude response highlights the differences in efficiency and ripple behavior among the filter types. The Butterworth filter exhibits a maximally flat passband, but it generally requires a higher order to meet the $60$ dB stopband requirement. In contrast, the Chebyshev Type I filter introduces passband ripples to achieve a steeper roll-off, while the Chebyshev Type II filter maintains a flat passband at the expense of ripples in the stopband. Among all the designs, the Elliptic filter provides the sharpest transition between the passband and stopband by allowing ripples in both regions, making it the most order-efficient option.

\begin{figure}[htbp]
    \centering
    \includegraphics[width=\columnwidth]{figure1.png}
\end{figure}

\subsection{Signal Reconstruction and the Gibbs Phenomenon}

The phase response, shown in Figure 2, illustrates the frequency-dependent phase shifts present in practical filter designs. While the ideal filter maintains a constant zero phase, all realizable analog filters exhibit some level of non-linearity. Among the filters, the Butterworth design generally provides the most linear phase response within the passband, which is important for preserving signal timing. However, as the filters become more aggressive in terms of magnitude roll-off—particularly in the Chebyshev and Elliptic designs—the phase distortion becomes more pronounced, especially near the cutoff frequencies of $2000$ and $4000$ rad/sec.

\begin{figure}[htbp]
    \centering
    \includegraphics[width=\columnwidth]{figure2.png}
\end{figure}